% §5 The cascade --- blueprint chapter 5 (blueprint/src/parts/05-cascade.tex), statements and
% proofs of record verbatim. Drafted 2026-09-11.
%
% Contents and departures from the blueprint chapter:
%   - all four statement nodes transcribed (lem:additivity, thm:increments-levy,
%     cor:monotonicity, cor:smoothed-transmittance) under their markers; all [T] and \leanok,
%     and the theorem's proof is the machine-checked route with its two constants (the sharp
%     ones survive inside the proof as a remark marked not machine-checked, as in the blueprint);
%   - the lead paragraph adds the draft's "three things" sentence; the paragraph before the
%     theorem is rewritten in half-line-versus-line terms, and the comparisons with the causal
%     article that the blueprint spreads over its annotations are collected in this section's
%     one relation-to-the-causal-theory remark, rem:causal-relation-cascade, after
%     cor:monotonicity (the proof of thm:increments-levy keeps its one verbatim clause);
%   - rem:lattice-family is copied; the probabilistic reading of Theta in
%     cor:smoothed-transmittance's annotation is kept as prose after the corollary.
% Forward references: lem:no-lattice, lem:action-rigidity, sec:covariance (§6),
% prop:strict-positivity (§7).

\section{The cascade: additivity and infinite divisibility}\label{sec:cascade}

With the representation in hand, the cascade axiom (A6) becomes a statement about kernels:
composition of operators is convolution of measures, so the exponent of
\S\ref{sec:representation} is additive over adjacent intervals of the scale axis. The optical
densities of stacked stages add (Remark~\ref{rem:modulation-transfer}). Set
\[
  G(t,\omega) := g_{0,t}(\omega),
\]
the density accumulated from scale $0$ to scale $t$. The exponent of any increment is a
difference of its values, so $G(\cdot,\omega)$ is nondecreasing; every increment has the
symmetric \Levy--Khintchine form \ref{eq:levy-khintchine}, which is the hemigroup form of
infinite divisibility; and, under nondegeneracy, $G(\cdot,\omega)$ is strictly increasing
for all but countably many $\omega$. The
coordinate that \S\ref{sec:covariance} uses on the scale axis is therefore not
$G(\cdot,\omega_0)$ but a smoothed transmittance, Corollary~\ref{cor:smoothed-transmittance},
which is strictly monotone with no exceptions. No covariance is used anywhere in this
section: everything below holds for every family satisfying (A1)--(A7), with (ND) added only
where a statement says so, including the lattice family of Remark~\ref{rem:lattice-family}.

% shared with blueprint lem:additivity
\begin{lemma}[additivity]\label{lem:additivity}
Under (A1)--(A7), $\mu_{r,t} = \mu_{s,t} * \mu_{r,s}$ for $r \le s \le t$, hence
\[
  g_{r,t}(\omega) = g_{r,s}(\omega) + g_{s,t}(\omega), \qquad g_{t,t} = 0 .
\]
Consequently $g_{s,t}(\omega) = G(t,\omega) - G(s,\omega)$, where $G(0,\cdot) = 0$;
$G(\cdot,\omega)$ is nondecreasing and continuous for each $\omega$; and $G(t,\cdot)$ is even
and continuous with $G(t,0) = 0$.
\end{lemma}

\begin{proof}
(A6) gives $\Phis{s}{t}\Phis{r}{s} = \Phis{r}{t}$, which by
Lemma~\ref{lem:convolution-representation} reads $\mu_{s,t} * (\mu_{r,s} * f) = \mu_{r,t} * f$ for
every $f \in \Lone$, so $\mu_{r,t} = \mu_{s,t} * \mu_{r,s}$ by the uniqueness clause there.
Fourier transforms turn convolution into multiplication and hence $-\log$ into addition, which
is legitimate because $\hat\mu > 0$ (Lemma~\ref{lem:nonvanishing}); and $\Phis{t}{t} = \mathrm{Id}$
gives $\mu_{t,t} = \delta_0$ and $g_{t,t} = 0$. The remaining clauses are
Lemma~\ref{lem:nonvanishing}, together with $g_{s,t} \ge 0$ for the monotonicity of
$G(\cdot,\omega)$.
\end{proof}

In the semigroup case the next step is trivial, $\mu_t = \mu_{t/n}^{*n}$ exhibiting every
kernel as a convolution power. In the hemigroup case there is no such power, and what replaces
it is a null-array limit. The idea is the following. Split the increment from $s$ to $t$
into $n$ consecutive stages. By additivity the exponent $g_{s,t}$ is the sum of the $n$
small exponents, and by continuity every one of them is close to zero once $n$ is large.
For a small exponent $g_i$ the kernel $\mu_i$ is close to a point mass at the origin, and
$g_i(\omega) \approx 1 - \hat\mu_i(\omega) = \int (1 - \cos\omega x)\,\mu_i(dx)$. Summing
over the stages, $g_{s,t}(\omega)$ is the limit of $\int (1 - \cos\omega x)\,\Pi_n(dx)$,
where $\Pi_n$ is the sum of the $n$ kernels. If the measures $\Pi_n$ converged to a measure
$\nu$, the limit would be of the form \ref{eq:levy-khintchine}. They do not: $\Pi_n$ has
total mass $n$, and that mass piles up at the origin. The remedy is to move the factor
$1 \wedge x^2$ from the integrand to the measure, writing
$\int (1 - \cos\omega x)\,\Pi_n(dx) = \int \frac{1 - \cos\omega x}{1 \wedge x^2}\,\varrho_n(dx)$
with $\varrho_n(dx) := (1 \wedge x^2)\,\Pi_n(dx)$. The new integrand is bounded, and near the
origin it tends to $\omega^2/2$, so the mass that piles up there produces the Gaussian term
$a\,\omega^2$ rather than a divergence. The measures $\varrho_n$ have bounded mass, and a
limit of them along a subsequence gives the pair $(a,\nu)$.

One obstacle remains. The new integrand oscillates as $x \to \infty$ instead of converging,
so it does not extend continuously to $x = \infty$, and mass of $\varrho_n$ escaping to
infinity would be lost from the limit rather than accounted for. So it has to be shown,
before the limit is taken, that no mass escapes: the family $\varrho_n$ is tight. The tool
is a truncation inequality, and its only input is that $g_{s,t}$ is continuous at
$\omega = 0$ with $g_{s,t}(0) = 0$.

% shared with blueprint thm:increments-levy
\begin{theorem}[increments are symmetric \Levy{} exponents]\label{thm:increments-levy}
Under (A1)--(A7), $g_{s,t} \in \LEs$ for every $s \le t$: there are a Gaussian coefficient
$a(s,t) \ge 0$ and a measure $\nu_{s,t}$ on $(0,\infty)$ with
$\int_{(0,\infty)}(1 \wedge x^2)\,\nu_{s,t}(dx) < \infty$ and
\begin{equation}\label{eq:increment-levy}
  g_{s,t}(\omega) \;=\; a(s,t)\,\omega^2
   + \int_{(0,\infty)} \bigl(1 - \cos\omega x\bigr)\,\nu_{s,t}(dx),
\end{equation}
the pair being unique by Proposition~\ref{prop:fourier-toolbox}(3). In particular every kernel
$\mu_{s,t}$ is infinitely divisible.
\end{theorem}

The representation is extracted directly, and the closure clause
Proposition~\ref{prop:fourier-toolbox}(4) is never used. What
Proposition~\ref{prop:fourier-toolbox}(3) supplies is uniqueness of the pair and, for the
last sentence, the converse direction. The null-array limit itself, including the
truncation inequality that excludes escape of mass, is proved below, and Prokhorov's theorem
is the one classical fact it stands on. It is a theorem of the library, on Lean core, and not a
cited interface. The proof is the machine-checked route: its two
constants, $u^2$ and $\frac{2}{3\pi^2}$, are those of the verified inequalities rather
than the sharp $u^2/2$ and $\frac{2}{15}$, which are recorded at the end of the second step
and are not machine-checked. In the Lean development the representation itself rests on
Lean core alone, and only the closing sentence uses two clauses of the toolbox,
Proposition~\ref{prop:fourier-toolbox}(1) and the converse half of (3).

\begin{proof}
Fix $s < t$ and the uniform partition $s_i = s + \tfrac in(t-s)$, $i = 0,\dots,n$; write
$\mu_i := \mu_{s_i,s_{i+1}}$ and $g_i := g_{s_i,s_{i+1}}$, so that
$g_{s,t} = \sum_{i<n} g_i$ for every $n$, by Lemma~\ref{lem:additivity}.

\emph{The null array.} Fix $\omega$. $G(\cdot,\omega)$ is uniformly continuous on $[s,t]$, so
$m_n := \max_i g_i(\omega) \to 0$. For $u \ge 0$ we have $0 \le u - (1 - e^{-u}) \le u^2$: the
left inequality is $e^v \ge 1 + v$ at $v = -u$, and the right one is the same inequality at
$v = u$ followed by $e^{-u} \le (1+u)^{-1} \le 1 - u + u^2$. Hence
\[
  \Bigl|\,g_{s,t}(\omega) - \sum_i \bigl(1 - e^{-g_i(\omega)}\bigr)\Bigr|
  \ \le\ \sum_i g_i(\omega)^2
  \ \le\ m_n\,g_{s,t}(\omega) \ \to\ 0 .
\]
On the other hand, with $\Pi_n := \sum_i \mu_i$, a symmetric finite measure of mass $n$,
\[
  \sum_i \bigl(1 - e^{-g_i(\omega)}\bigr) = \sum_i \bigl(1 - \hat\mu_i(\omega)\bigr)
  = \int_\RR \bigl(1 - \cos\omega x\bigr)\,\Pi_n(dx)
  = \int_{(0,\infty)} \bigl(1 - \cos\omega x\bigr)\,\widetilde\Pi_n(dx),
\]
where $\widetilde\Pi_n$ is the image of $\Pi_n$ under $x \mapsto |x|$, the integrand being
even; the origin carries no weight, since $1 - \cos 0 = 0$, so $\widetilde\Pi_n$ may be taken
on $(0,\infty)$.

\emph{Two bounds, uniform in $n$.} Both come from one inequality. For $\eta > 0$, Tonelli and
the identity $\eta^{-1}\int_0^\eta (1 - \cos\omega x)\,d\omega = 1 - \frac{\sin\eta x}{\eta x}$
give
\begin{equation}\label{eq:truncation}
  \int_{(0,\infty)} \Bigl(1 - \frac{\sin\eta x}{\eta x}\Bigr)\,\widetilde\Pi_n(dx)
  \;=\; \frac1\eta\int_0^\eta \sum_i\bigl(1 - \hat\mu_i(\omega)\bigr)\,d\omega
  \;\le\; \frac1\eta\int_0^\eta g_{s,t}(\omega)\,d\omega \;=:\; \bar g(\eta),
\end{equation}
using $1 - e^{-u} \le u$ and Lemma~\ref{lem:additivity}. The function $1 - \frac{\sin u}{u}$ is
nonnegative and satisfies
\begin{equation}\label{eq:sinc-lower}
  1 - \frac{\sin u}{u} \ \ge\ \frac{2}{3\pi^2}\,\bigl(1 \wedge u^2\bigr)
  \qquad (u \in \RR),
\end{equation}
with the convention $\frac{\sin 0}{0} = 1$. Both sides being even, take $u > 0$. For
$0 < u \le \pi$, integrate the elementary bound $1 - \cos v \ge \frac{2}{\pi^2}v^2$, valid for
$|v| \le \pi$, over $[0,u]$ against the identity $u - \sin u = \int_0^u (1 - \cos v)\,dv$; this
gives $u - \sin u \ge \frac{2}{3\pi^2}u^3$, that is
$1 - \frac{\sin u}{u} \ge \frac{2}{3\pi^2}u^2 \ge \frac{2}{3\pi^2}(1 \wedge u^2)$. For
$u > \pi$ the crude bound $\frac{\sin u}{u} \le \frac1u < \frac1\pi$ leaves
$1 - \frac{\sin u}{u} > 1 - \frac1\pi > \frac23 > \frac{2}{3\pi^2} \ge
\frac{2}{3\pi^2}(1 \wedge u^2)$. \emph{Mass:} with $\eta = 1$, \ref{eq:truncation} and
\ref{eq:sinc-lower} bound the weighted measures
$\varrho_n := (1 \wedge x^2)\,\widetilde\Pi_n$ on $(0,\infty)$ by
$\varrho_n((0,\infty)) \le \tfrac{3\pi^2}{2}\,\bar g(1)$. \emph{Tightness:} for $x \ge 1/\eta$
the argument $u = \eta x$ has $1 \wedge u^2 = 1$, so \ref{eq:truncation} and
\ref{eq:sinc-lower} give
$\tfrac{2}{3\pi^2}\,\widetilde\Pi_n([1/\eta,\infty)) \le \bar g(\eta)$, and
$\bar g(\eta) \to 0$ as $\eta \downarrow 0$ because $g_{s,t}$ is continuous with
$g_{s,t}(0) = 0$ (Lemma~\ref{lem:nonvanishing}); since $\varrho_n \le \widetilde\Pi_n$ on
$[1,\infty)$, the tails of $\varrho_n$ are small uniformly in $n$. This is the truncation
inequality of characteristic-function theory, proved here in the few lines it takes; it is
where continuity at $\omega = 0$ does quantitative work.

Sharper constants are available and are not used: Taylor's theorem gives
$u - (1 - e^{-u}) \le u^2/2$ and, through $\sin u \le u - u^3/6 + u^5/120$,
$1 - \frac{\sin u}{u} \ge \frac{2}{15}(1 \wedge u^2)$. Neither is worth its derivative
argument here --- the first constant multiplies a quantity that is going to zero, the second
one that has only to be finite --- and neither is machine-checked.

\emph{The limit.} A family of finite measures on $[0,\infty)$ of uniformly bounded mass and
uniformly small tails has a weakly convergent subsequence: pass first to a subsequence along
which the total masses converge; if their limit is $0$ the measures converge weakly to the
zero measure, and otherwise Prokhorov's theorem applies to the normalised measures, which
are tight. Let $\varrho_{n_k} \to \varrho$, a finite measure on $[0,\infty)$, the subsequence
chosen once and independently of $\omega$. The function
\[
  k_\omega(x) := \frac{1 - \cos\omega x}{1 \wedge x^2}, \qquad k_\omega(0) := \tfrac12\omega^2,
\]
is continuous on $[0,\infty)$ --- the singularity is removable, since
$1 - \cos u = \tfrac12u^2 + O(u^4)$ --- and bounded by $\max(\tfrac12\omega^2, 2)$, so
$\int k_\omega\,d\varrho_{n_k} \to \int k_\omega\,d\varrho$, while
$\int k_\omega\,d\varrho_n = \int (1-\cos\omega x)\,\widetilde\Pi_n(dx) \to g_{s,t}(\omega)$ by
the first step. Reading $\varrho$ in two pieces, the atom at $0$ and the rest,
\[
  g_{s,t}(\omega) = \tfrac12\varrho(\{0\})\,\omega^2
   + \int_{(0,\infty)} \bigl(1 - \cos\omega x\bigr)\,\frac{\varrho(dx)}{1 \wedge x^2},
\]
which is \ref{eq:increment-levy} with $a(s,t) := \tfrac12\varrho(\{0\})$ and
$\nu_{s,t}(dx) := (1 \wedge x^2)^{-1}\varrho(dx)$ on $(0,\infty)$, for which
$\int (1 \wedge x^2)\,\nu_{s,t} = \varrho((0,\infty)) < \infty$. The atom at $0$ is the Gaussian
coefficient, as in the causal theory it was the drift; there is no atom at infinity to exclude,
because tightness has already forbidden it, and correspondingly no killing term. The case
$s = t$ is $g_{t,t} = 0$, with the pair $(0,0)$.

\emph{Infinite divisibility.} For each $n$, $g_{s,t}/n$ is of the form
\ref{eq:levy-khintchine} with pair $(a/n, \nu/n)$, so by the converse half of
Proposition~\ref{prop:fourier-toolbox}(3) and by Proposition~\ref{prop:fourier-toolbox}(1),
$e^{-g_{s,t}/n}$ is the transform of a symmetric probability measure, whose $n$-fold
convolution has transform $e^{-g_{s,t}} = \hat\mu_{s,t}$ and hence is $\mu_{s,t}$ by
Proposition~\ref{prop:fourier-uniqueness}.
\end{proof}

Note the order of dependence. Infinite divisibility implies a zero-free transform, so
Lemma~\ref{lem:nonvanishing} is now also a corollary of the theorem. But the theorem's proof
used the lemma to define the exponents it manipulates, and the order cannot be reversed.

% shared with blueprint cor:monotonicity
\begin{corollary}[monotonicity, and where it is strict]\label{cor:monotonicity}
Assume in addition (ND). For $s < t$ the zero set
$N_{s,t} := \{\omega : g_{s,t}(\omega) = 0\}$ is a closed proper subgroup of $\RR$, hence
$\{0\}$ or a lattice $c\ZZ$ with $c > 0$, and it is a lattice exactly when $\mu_{s,t}$ is a
lattice law. Consequently $G(\cdot,\omega)$ is strictly increasing on $[0,\infty)$ for every
$\omega$ outside the countable set $E := \bigcup N_{p,q}$, the union over rational
$0 \le p < q$.
\end{corollary}

\begin{proof}
$N_{s,t} = N(\mu_{s,t})$ in the notation of Lemma~\ref{lem:lattice-zero}, a closed subgroup,
proper because $\mu_{s,t} \ne \delta_0$ by (ND) and the uniqueness clause of
Lemma~\ref{lem:convolution-representation}. If $[p,q] \subseteq [s,t]$ then
$g_{s,t} = g_{s,p} + g_{p,q} + g_{q,t}$ with all three terms nonnegative
(Lemmas~\ref{lem:additivity} and~\ref{lem:nonvanishing}), so $N_{s,t} \subseteq N_{p,q}$;
every interval $[s,t]$ with $s < t$ contains one with rational endpoints, so every $N_{s,t}$
is contained in some $N_{p,q}$ with $p < q$ rational, and $E$ is a countable union of
lattices, hence countable, and it contains $0$, every $N_{p,q}$ being a subgroup. For
$\omega \notin E$ and $s < t$ we get $g_{s,t}(\omega) > 0$, and
$G(t,\omega) - G(s,\omega) = g_{s,t}(\omega)$ is Lemma~\ref{lem:additivity}.
\end{proof}

This is as much as the line gives before covariance: the exceptional frequencies are exactly
those at which some increment is a lattice law, and strict monotonicity at every frequency
becomes true again only after the classification, Proposition~\ref{prop:strict-positivity}(2).
The exceptional set can be larger than $\{0\}$.

% blueprint rem:lattice-family, copied
\begin{remark}[the exceptional set is not $\{0\}$]\label{rem:lattice-family}
For the symmetric compound Poisson family on $\ZZ$ with
$\hat\mu_{s,t}(\omega) = e^{-(t-s)(1-\cos\omega)}$ (a hemigroup, indeed a semigroup, of unit
steps), every axiom except (A8) holds, every kernel is a lattice law,
$E = 2\pi\ZZ$, and $G(\cdot,2\pi n) \equiv 0$ for every integer $n$. Covariance
fails because $\Dil{\lambda}\mu_{s,t}$ has exponent $(t-s)(1-\cos\lambda\omega)$, which is of
the form $c\,(1-\cos\omega)$ only for $\lambda = 1$. So the exclusion of lattice kernels is a
consequence of covariance rather than a property of the function class. That is
Lemma~\ref{lem:no-lattice}.
\end{remark}

\begin{remark}[relation to the causal theory]\label{rem:causal-relation-cascade}
The first two statements of this section port from the causal article
\sscite{fagerstrom2026hemigroup}, the additivity lemma verbatim and the increments theorem
with the weight $1 - e^{-x}$ replaced by $1 \wedge x^2$ and the compactification of the
half-line replaced by the truncation inequality \ref{eq:truncation}, which is the one new
idea. The atom at the origin that was the drift there is the Gaussian coefficient here, and
the atom at infinity that had to be excluded there cannot form here, because tightness has been
established first. The third statement does not port verbatim.
The causal corollary says that the accumulated exponent is injective in scale at every
value of the transform variable, and that this is what makes the relabelling of the scale
axis rigid. On the line that sentence is false before covariance,
Remark~\ref{rem:lattice-family} being the counterexample, and
Corollary~\ref{cor:monotonicity} states what survives. The rigidity argument of
\S\ref{sec:covariance} is therefore run on a different coordinate, the smoothed
transmittance below, which has no causal counterpart because none was needed.
\end{remark}

The coordinate is one number attached to each scale, strictly monotone in it, and obtained
without covariance. It is the transform of the kernel averaged against a window on the
frequency axis. Any strictly positive integrable window would serve: the average is
strictly monotone because a nondegenerate increment cannot have transform $1$ on a set of
positive measure, whereas a single frequency can fall in the zero set of a lattice
increment. The Gaussian is chosen because it is its own transform and because its dilates
are dominated by one Gaussian, which is what the continuity estimate of
\S\ref{sec:covariance} uses.

% shared with blueprint cor:smoothed-transmittance
\begin{corollary}[the smoothed transmittance]\label{cor:smoothed-transmittance}
Assume (A1)--(A7) and let $\rho$ be the standard Gaussian density. Put
\[
  \Theta(t) := \int_\RR \hat\mu_{0,t}(\omega)\,\rho(\omega)\,d\omega
   \;=\; \int_\RR e^{-x^2/2}\,\mu_{0,t}(dx) \ \in\ (0,1].
\]
Then $\Theta$ is continuous on $[0,\infty)$, $\Theta(0) = 1$, and $\Theta$ is nonincreasing;
under (ND) it is strictly decreasing.
\end{corollary}

\begin{proof}
The two expressions agree by Fubini, since $\int e^{-i\omega x}\rho(\omega)\,d\omega =
e^{-x^2/2}$ exactly for $\rho(\omega) = (2\pi)^{-1/2}e^{-\omega^2/2}$. Continuity in $t$ is
dominated convergence, with the bound $|\hat\mu| \le 1$ and the continuity of
$\hat\mu_{0,t}(\omega)$ in $t$ (Lemma~\ref{lem:nonvanishing}); positivity is that
$e^{-x^2/2} > 0$ at every $x$ and $\mu_{0,t}$ is a probability measure, and
$\Theta(0) = 1$ is $\mu_{0,0} = \delta_0$. For $s < t$, $\hat\mu_{0,t} =
\hat\mu_{s,t}\,\hat\mu_{0,s}$ (Lemma~\ref{lem:additivity}) with $0 < \hat\mu_{s,t} \le 1$, so
\[
  \Theta(t) - \Theta(s)
  = \int_\RR \bigl(\hat\mu_{s,t} - 1\bigr)\,\hat\mu_{0,s}\,\rho\,d\omega \ \le\ 0 ;
\]
under (ND) the continuous integrand is negative on the open set $\{\hat\mu_{s,t} < 1\}$, which
is nonempty because $\mu_{s,t} \ne \delta_0$ (Lemma~\ref{lem:convolution-representation}),
while $\hat\mu_{0,s}\,\rho > 0$, so the inequality is strict.
\end{proof}

Probabilistically $\Theta(t) = \Ex\,e^{-X_t^2/2}$ with $X_t := X_{0,t}$: the transmittance of
the kernel from the origin averaged over a Gaussian spectrum of probe wavenumbers, or equally the
probability that $X_t^2/2$ falls short of an independent standard exponential threshold.
It is what Lemma~\ref{lem:action-rigidity}(3) inverts.
Note that the positivity of $\Theta$ is not Lemma~\ref{lem:nonvanishing} but the positivity
of the integrand $e^{-x^2/2}$. The positivity of $\hat\mu$ is used only on the two
monotonicity clauses.
